\subsection{Summary of Contributions}

This thesis has presented QASAMAP (the Quantum-Agentic Smart Manufacturing Anomaly detection and Predictive Maintenance) framework as a principled, integrated solution to the multidimensional challenge of intelligent maintenance decision support in modern industrial environments. The research was grounded in a 100,000-reading smart manufacturing dataset representing 50 machines monitored across 69 days, whose statistical properties an 8.92\% anomaly rate, five distinct failure modalities, 16,333 critical low-RUL readings, and temporal non-stationarity provided both the motivation and the empirical basis for every architectural decision in QASAMAP. Building from this empirical foundation, the thesis conducted systematic literature synthesis, gap analysis, architectural design, comprehensive experimental evaluation, and critical assessment of limitations.

The \textbf{first primary contribution} is the T-GCN spatiotemporal forecasting architecture with cross-correlation lag graph construction. By building a directed graph over 50 machines (50 nodes, 538 edges) from pairwise lagged cross-correlations of the temperature time series, and combining Graph Convolution layers for spatial aggregation with GRU layers for temporal modeling, the T-GCN achieves an average R\textsuperscript{2} of 0.7702 across five sensor channels. This represents a 12.4\% improvement over the OptimizedLSTM baseline (R\textsuperscript{2} = 0.6853) and an 8.0\% improvement over the Hybrid LSTM-GCN (R\textsuperscript{2} = 0.7134), demonstrating that graph-based cross-machine dependency modeling provides a meaningful and consistent benefit for industrial sensor stream forecasting.

The \textbf{second primary contribution} is a comprehensive empirical benchmark of five supervised learning algorithms (XGBoost, LightGBM, CatBoost, HistGradientBoosting, TabNet) across six machine health prediction targets. Tree-based ensemble models achieve near-perfect performance on binary classification targets such as anomaly flag (best F1 = 0.999, LightGBM) and downtime risk (best F1 = 0.999, LightGBM), reflecting the high discriminability of these targets. More challenging tasks --- machine status (best F1 = 0.712, XGBoost), failure type (best F1 = 0.880, XGBoost), and maintenance required (best F1 = 0.872, XGBoost) --- provide meaningful performance benchmarks. The RUL regression results (best R\textsuperscript{2} = 0.189, CatBoost) confirm the fundamental difficulty of single-snapshot RUL prediction without temporal degradation trajectory context, motivating the use of the T-GCN architecture for this task.

The \textbf{third primary contribution} is the Pauli-Z multi-sensor risk encoding and quantum optimization pipeline. The \texttt{PauliZRiskEncoder} maps each machine's five normalized sensor values to a composite Pauli-Z expectation $\langle Z \rangle_c$, enabling detection of sub-threshold risk combinations that classical single-sensor thresholds miss. Combined with QUBO-based Quantum Annealing for machine priority ranking and QAOA for maintenance slot scheduling, the quantum layer triggers actionable overrides for 5 out of 7 machines in the experimental evaluation machines that the classical risk engine would route to passive monitoring. This demonstrates that quantum sensing provides a complementary and practically valuable signal to classical threshold-based reasoning.

The \textbf{fourth primary contribution} is the multi-agent agentic AI orchestration system, comprising Monitoring, Diagnosis, Planning, and Reporting agents coordinated by the \texttt{MultiAgentOrchestrator} via the \texttt{gpt-oss:120b} model on Ollama Cloud. Supported by the \texttt{ContextMemoryStore}, \texttt{ReasoningEngine}, quantum override logic, and \texttt{LearningLoop}, the system autonomously manages the complete predictive maintenance decision cycle from sensor event ingestion through maintenance action dispatch. The three-way comparative evaluation demonstrates that the full QASAMAP stack achieves detection accuracy of 94.1\%, false negative rate of 4.7\%, mean response time of 3.2 minutes, machine coverage of 85.7\%, 58.3\% downtime reduction, and 43.9\% maintenance cost savings improvements of $+$53.8\%, $-$84.9\%, $-$92.9\%, $+\infty$, $+\infty$, and $+\infty$ respectively over the traditional rule-based baseline.

\subsection{Research Objective Verification}

The five research objectives formulated in Chapter 1 have been addressed and the following conclusions are drawn.

\textbf{Objective 1} (T-GCN spatiotemporal forecasting) is fully achieved. The T-GCN architecture with cross-correlation lag graph construction, machine-wide propagation with anomaly injection, and real-time streaming with partial sensor imputation via \texttt{AnchorBufferManager} is implemented, evaluated, and demonstrated to outperform both LSTM and Hybrid LSTM-GCN baselines on the 100,000-reading dataset (avg.\ R\textsuperscript{2} = 0.7702 vs.\ 0.6853 and 0.7134).

\textbf{Objective 2} (Multi-target supervised classification and regression) is fully achieved. Five classification targets and one regression target are evaluated across five algorithms, providing complete performance benchmarks for all machine health prediction tasks defined in the dataset schema.

\textbf{Objective 3} (Quantum computing integration) is substantially achieved. Pauli-Z multi-sensor encoding, QUBO-based quantum annealing, and QAOA scheduling are implemented using Qiskit Aer simulation and demonstrated on the 7-machine experimental set, with quantum overrides triggering actionable decisions for 5 of 7 machines. The current implementation uses classical simulation rather than physical quantum hardware, which constitutes the primary caveat on this objective.

\textbf{Objective 4} (Multi-agent agentic AI system) is fully achieved. The four-agent pipeline (Monitoring, Diagnosis, Planning, Reporting) with quantum override integration, contextual memory, and continuous learning is implemented and demonstrated to produce interpretable, quantum-informed maintenance decisions for all machines in the experimental machine set.

\textbf{Objective 5} (Three-way comparative evaluation) is fully achieved. The QASAMAP framework is benchmarked against a Traditional rule-based system and an Agentic AI Only configuration, with eight KPIs quantified across all three approaches and the independent contributions of the agentic and quantum layers both identified and quantified.

\subsection{Limitations and Critical Reflections}

Scientific integrity demands honest assessment of the limitations of this research alongside its contributions. The most significant limitation is the reliance on classical simulation of quantum circuits rather than execution on physical quantum hardware. All quantum computing results in this thesis are obtained from Qiskit Aer statevector simulations that assume perfect qubit fidelity, zero gate error rates, and idealized noise models. The transition to physical NISQ hardware will require addressing gate error rates, qubit decoherence, and limited connectivity that incur additional circuit overhead. The quantum computing results should therefore be interpreted as upper bounds on achievable quantum advantage under ideal conditions, pending validation on physical hardware.

The experimental evaluation of the quantum-agentic pipeline used a representative sample of 11 sensor events across 7 machines, which is sufficient to demonstrate the decision-making behavior of the full pipeline but does not constitute a large-scale statistical evaluation over the full 100,000-observation dataset. The KPI metrics reported for the three-way comparison (detection accuracy, false positive/negative rates, response time, etc.) are derived from the comparative analysis implementation in the codebase and represent scenario-level performance characterizations rather than cross-validated estimates over independent test sets. Future work should conduct larger-scale evaluation over the full dataset with cross-validation.

The supervised classification benchmarks reveal that the anomaly flag and downtime risk targets are near-perfectly discriminable by all tested models (F1 $> 0.998$), which suggests these labels may be directly derivable from deterministic rules rather than representing a genuine learning challenge. The more practically meaningful classification targets are machine status, failure type, and maintenance required, where F1 scores in the range 0.71--0.88 reflect genuine learning difficulty. The RUL regression results (R\textsuperscript{2} $\leq 0.19$) across all models indicate that static feature-based RUL prediction is fundamentally limited, reinforcing the importance of the T-GCN temporal approach for this task.

The dataset, while large at 100,000 observations, represents a single industrial context and may not generalize universally across manufacturing environments, machine types, or failure mode distributions. Validation on publicly available benchmark datasets such as the NASA C-MAPSS turbofan degradation dataset, the PRONOSTIA bearing dataset, and the IMS bearing dataset would strengthen the generalizability of the findings.

\subsection{Future Research Directions}

The most immediately impactful future direction is the evaluation of the QASAMAP quantum computing components on physical quantum hardware. The IBM Quantum Network, Google Quantum AI, and IonQ platforms offer processors with 27--127 qubits that approach the specifications required for the quantum circuits implemented in this thesis. Executing the Pauli-Z encoding circuits, QUBO annealing, and QAOA scheduler on physical hardware, characterizing noise-induced performance degradation, and developing hardware-specific error mitigation strategies would constitute a direct and important extension of the present work.

The second high-priority direction is the large-scale cross-validated evaluation of the full QASAMAP pipeline over the complete 100,000-observation dataset, replacing the scenario-based pipeline demonstration with rigorous statistical benchmarking over held-out test periods. This would provide statistical confidence intervals on all reported KPI metrics and enable direct comparison with the supervised learning baselines established in the classification and regression experiments.

The third direction is the extension of the multi-agent agentic layer to support enhanced LLM-based reasoning capabilities, including in-context learning from historical decision logs, chain-of-thought diagnostic reasoning over multi-step sensor degradation trajectories, and autonomous generation of detailed maintenance procedure reports. This would enhance the practical utility of the Reporting Agent output and reduce the domain expertise required to interpret and act on the system's recommendations.

The fourth direction is the investigation of federated learning extensions that enable multiple manufacturing facilities to collaboratively train shared predictive maintenance models without sharing raw sensor data. The multi-agent architecture of QASAMAP is naturally compatible with federated paradigms, and cross-facility collaboration would substantially increase the diversity and volume of training data available for rare failure mode detection.

\subsection{Broader Implications}

The research contributions of this thesis have implications extending beyond smart manufacturing predictive maintenance. The demonstration that Pauli-Z quantum encoding combined with QUBO annealing and QAOA scheduling produces actionable machine prioritization and slot assignment decisions in a real-world industrial context advances the practical case for quantum-enhanced industrial intelligence, even in the current simulation-based phase of quantum computing development. As quantum hardware continues to mature, the pipeline architectures demonstrated in this thesis provide a concrete template for transitioning quantum-enhanced predictive maintenance from research prototype to production deployment.

The demonstration that a four-agent LLM orchestration system with quantum override integration can autonomously manage the complete predictive maintenance decision cycle from sensor event ingestion through maintenance action dispatch, feedback collection, and threshold adaptation illustrates a practical model for human-AI collaboration in industrial operations. Rather than positioning AI as a passive analytical tool, the QASAMAP framework establishes a model of active AI agency in which automated agents handle continuous monitoring, diagnostic reasoning, and routine action dispatch while human supervisors retain oversight authority over high-risk decisions through the human-in-the-loop approval gate.

\subsection{Closing Remarks}

This thesis has demonstrated that the integration of spatiotemporal deep learning, quantum computing, and multi-agent agentic AI within a unified, closed-loop predictive maintenance framework is both technically feasible and empirically beneficial. The QASAMAP framework, grounded in the concrete evidence of a real-world industrial dataset, advances the state of the art across three intersecting domains while maintaining focus on the operational requirements of accuracy, adaptability, interpretability, and autonomous decision-making. The T-GCN forecasting, supervised learning benchmarks, quantum-agentic pipeline, and three-way comparative evaluation collectively constitute a coherent contribution to the emerging interdisciplinary area of quantum-agentic industrial intelligence. As quantum hardware matures and agentic AI capabilities expand, the architectural principles and empirical findings of this thesis provide a productive foundation for advancing the frontier of intelligent manufacturing.